\documentclass[11pt,a4paper]{article}

\usepackage[utf8]{inputenc}
\usepackage[T1]{fontenc}
\usepackage[french]{babel}
\usepackage{geometry}
\geometry{margin=2.5cm}
\usepackage{enumitem}
\usepackage{hyperref}
\usepackage{fancyhdr}

\pagestyle{fancy}
\fancyhf{}
\rhead{Cahier des Charges}
\lhead{Architecture Réseau Multi-Sites}
\rfoot{\thepage}

\title{\textbf{Cahier des Charges}\\
	Conception et Mise en Œuvre d’une Architecture Réseau Multi-Sites avec Extension Cloud}
\author{Nozha Safi (Jlidat)}
\date{\today}

\begin{document}
	
	\maketitle
	\thispagestyle{empty}
	\newpage
	\tableofcontents
	\newpage
	
	%-------------------------------------------------
	\section{Préambule}
	
	Le présent document définit les exigences techniques relatives à la conception
	et à la mise en œuvre d’une architecture réseau multi-sites intégrant une
	extension vers un fournisseur de Cloud.
	
	L’objectif est de structurer l’infrastructure interne, d’assurer la
	segmentation logique des départements, de sécuriser l’accès Internet et
	de permettre l’hébergement hybride de services applicatifs.
	
	%-------------------------------------------------
	\section{Contexte}
	
	L’entreprise dispose de deux sites géographiquement éloignés :
	
	\begin{itemize}
		\item Site 1
		\item Site 2
	\end{itemize}
	
	Chaque site héberge :
	
	\begin{itemize}
		\item Un département Administratif
		\item Un département Développeurs
		\item Des serveurs Intranet
		\item Des serveurs accessibles depuis Internet
	\end{itemize}
	
	Afin d’absorber une croissance d’activité, l’entreprise externalise
	une partie de ses serveurs auprès d’un Fournisseur de Cloud.
	
	Dans un objectif de modernisation, de performance et de sécurité, l'entreprise souhaite :
	\begin{itemize}
		\item{$\star$} Définir un plan d’adressage IPv4 cohérent en respectant les contraintes de capacité.  
		\item{$\star$} Segmenter les différents départements via VLAN 802.1Q.   
		\item{$\star$} Isoler les serveurs internes des serveurs exposés  
		\item{$\star$} Contrôler les communications inter-réseaux
		\item{$\star$} Sécuriser l'accès Internet
		\item{$\star$} Intégrer une solution Cloud privée et publique 
	\end{itemize}
	%-------------------------------------------------
	
	\section{Objectifs}
	Le projet vise à :
	\begin{itemize}
		\item[$\rightarrow$] Concevoir une architecture réseau multi-sites cohérente
		\item[$\rightarrow$] Mettre en oeuvre des VLAN conformes à la normes IEEE 802.1Q
		\item[$\rightarrow$] Assurer le routage inter-VLAN
		\item[$\rightarrow$] Implémenter des règles de filtrage via ACL étendues
		\item[$\rightarrow$] Mettre en oeuvre des mécanismes NAT (statique et dynamique)
		\item[$\rightarrow$] Garantir la séparation logique des flux internes et extrnes
	\end{itemize}
	%-------------------------------------------------
	\section{Architecture Générale}
	
	Chaque site comprend :
	
	\begin{itemize}[label=$\bullet$]
		\item Un routeur d’accès Internet (RE1 pour Site 1, RE2 pour Site 2)
		\item Deux commutateurs interconnectés en trunk 802.1Q
		\begin{itemize}
			\item Un commutateur dédié aux utilisateurs
			\item Un commutateur dédié aux serveurs
		\end{itemize}
		\item Quatre sous-réseaux distincts
	\end{itemize}
	
	Chez le fournisseur Cloud :
	
	\begin{itemize}[label=$\bullet$]
		\item Cloud privé => Sous-réseau LAN\_3-31
		\item Cloud public => Sous-réseau LAN\_3-32
		\item Routeurs fournisseurs assurant l'interconnexion
		\begin{itemize}
			\item Routeur RF1 (réseau Intranet Cloud)
			\item Routeur RF2 (réseau Public Cloud)
		\end{itemize}
	\end{itemize}
	
	%-------------------------------------------------
	\section{Plan d’Adressage}
	
	\subsection{Site 1}
	
	Plage publique attribuée : 123.1.1.160/27
	
	Sous-réseaux à définir :
	
	\begin{itemize}
		\item LAN\_1-1 : 15 (14) hôtes (Administratif)
		\item LAN\_1-2 : 20 (18) hôtes (Développeurs)
		\item LAN\_1-31 : 7 (6) hôtes (Serveurs Intranet)
		\item LAN\_1-32 : jusqu’à 15 (14) adresses publiques
	\end{itemize}
	
	RE1 (côté SW2): IP (Broadcast : ) !!!!!
	
	Réseau privé : 192.168.1.0/24
	
	
	
	\subsection{Site 2}
	
	Plage publique attribuée : 123.2.2.160/27
	
	Sous-réseaux à définir :
	
	\begin{itemize}
		\item LAN\_2-1 : 10 (7) hôtes
		\item LAN\_2-2 : 25 (26) hôtes
		\item LAN\_2-31 : 4 (2) hôtes
		\item LAN\_2-32 : jusqu’à 7 (6) adresses publiques
	\end{itemize}
	
	RE2 (côté SW4) : IP (Broadcast : ) !!!!!
	
	Réseau privé :  192.168.2.0/24
	
	\subsection{Fournisseur de Cloud}
	
	Plage publique Cloud : 163.173.19.224/28
	
	Sous-réseaux à définir:
	
	\begin{itemize}
		\item LAN\_3-31 : 10 serveurs Intranet
		\item LAN\_3-32 : 10 serveurs Public
		\item Liaisons point-à-point :
		\begin{itemize}
			\item RF1 : 163.173.18.254/31
			\item RF2 : 163.173.19.222/31
		\end{itemize}
	\end{itemize}
	
	RF1 (côté SW5) : IP (Broadcast : ) !!!!!
	RF2 (côté SW6) : IP (Broadcast : ) !!!!!
	
	Réseau privé :  192.168.3.0/24
	
	%-------------------------------------------------
	\section{Segmentation VLAN}
	
	Chaque site met en œuvre des VLAN 802.1Q :
	
	\begin{itemize}
		\item VLAN 10 : Administratif
		\item VLAN 20 : Développeurs
		\item VLAN 31 : Serveurs Intranet
		\item VLAN 32 : Serveurs DMZ
	\end{itemize}
	
	Le routage inter-VLAN est assuré par RE1 et RE2 via encapsulation IEEE 802.1Q.
	
	%-------------------------------------------------
	\section{Politique de Sécurité}
	
	Les règles suivantes doivent être appliquées :
	
	\begin{itemize}
		\item Interdiction de communication entre VLAN 10 et VLAN 20
		\item Interdiction de communication entre VLAN 31 et VLAN 32
		\item VLAN 31 accessible uniquement depuis les VLAN internes
		\item VLAN 32 accessible depuis Internet et depuis les VLAN internes
		\item Accès Internet autorisé pour tous les VLAN (flux sortants uniquement) ; Tous les VLAN peuvent initier des connexions vers Internet.		
	\end{itemize}
	
	Le système devra permettre l’isolation des flux inter-VLAN.
	
	%-------------------------------------------------
	\section{Services Réseau}
	
	\subsection{DHCP}
	
	Un serveur DHCP est déployé sur chaque site pour attribuer
	dynamiquement les adresses IP des sous-réseaux privés.
	
	Les serveurs applicatifs disposent d’adresses IP statiques.
	
	\subsection{NAT}
	
	Les routeurs RE1 et RE2 implémentent :
	
	\begin{itemize}
		\item NAT dynamique (PAT) pour les postes internes
		\item NAT statique pour les serveurs DMZ
	\end{itemize}
	
	Une adresse publique est réservée pour l’accès Internet au serveur Web.
	
	%-------------------------------------------------
	\section{Livrables}
	
	\begin{itemize}
		\item Plan d’adressage détaillé
		\item Configurations VLAN
		\item Configurations DHCP
		\item Règles ACL
		\item Configuration NAT
		\item Simulation complète sous GNS3
	\end{itemize}
	
\end{document}
